% Generated from locale/tr/content/set-theory/cardinals/cp.tex; do not hand-edit.


\olfileid[tr]{sth}{cardinals}{cp}
\olsection{Cantor İlkesi}

\olref[ordinals][vn]{sec} kesimine dönün. İyi sıralanmış kümeleri
tartışıyor ve iyi sıralamaları temsil edecek nesneler bulunmasının iyi
olacağını öneriyorduk. Bunu göz önünde tutarak sıra sayılarını tanıttık;
sonra \olref[ordinals][ordtype]{ordtypesworklikeyouwant} içinde bunların
istediğimiz gibi davrandığını gösterdik:
\[
	\ordtype{A, <} = \ordtype{B, \lessdot}
	\text{ ancak ve ancak } \tuple{A, <} \isomorphic \tuple{B, \lessdot}.
\]
Daha da geriye, \olref[sfr][siz][equ]{sec} kesimine dönün. Orada naif
biçimde çalışarak bir kümenin ``büyüklüğü'' kavramını tanıttık. Özellikle
iki kümenin, ancak ve ancak bir !!a{bijection} $f\colon A\to B$ varsa
eş güçlü, $\cardeq{A}{B}$, olduğunu söyledik. Bu, iyi sıralama kavramından
özünde daha basittir: yalnızca !!{bijection}s ilgileniriz; sıra
izomorfizmlerindeki gibi !!{bijection}s ``herhangi bir yapıyı koruyup
korumadığı'' bizi ilgilendirmez.

Bütün bunlar açık bir düşünceye yol açar. İyi sıralamaları ölçmek için
belirli nesneler, \emph{sıra sayıları}, tanıttığımız gibi büyüklüğü ölçmek
için belirli nesneler, \emph{kardinal sayılar}, tanıtabiliriz. Bu bölümün
amacı budur.

Bu kardinal sayıların ne olacağını söylemeden önce sağlamaları gereken bir
ilke ortaya koymalıyız. $X$ kümesinin kardinalitesi için $\card{X}$ yazarak
şu koşulu sağlamalarını isteriz:
\[
	\card{A} = \card{B} \text{ ancak ve ancak } \cardeq{A}{B}.
\]
Buna \emph{Cantor İlkesi} diyeceğiz; çünkü bu ilkeyi açık biçimde düşünen
ilk kişi muhtemelen Cantor'du. (\emph{Hume İlkesi} ile ilişkisi hakkında
\olref[hp]{sec} içinde daha fazlasını söyleyeceğiz.) Öyleyse amacımız her
$X$ için $\card{X}$'i Cantor İlkesi'ni verecek biçimde tanımlamaktır.

